% ID-5-02 (TO) · Mathematical Foundations of Diffusion Models · IIT Jammu
% HW 1 — Random Variables, Joints, Conditionals, Bayes
% Build:  latexmk -pdf hw1.tex     (or: pdflatex hw1.tex)
\documentclass[11pt,a4paper]{article}

\usepackage[margin=2.4cm]{geometry}
\usepackage{amsmath,amssymb}
\usepackage{booktabs}
\usepackage[colorlinks=true,linkcolor=black,urlcolor=blue]{hyperref}

\newcommand{\pts}[1]{\hfill\textbf{(#1pts)}}
\newcommand{\pt}[1]{\hfill\textbf{(#1pt)}}
\newcommand{\total}[1]{\par\medskip\noindent\textbf{Total: #1 points}}
\newcommand{\PP}{\mathbb{P}}
\newcommand{\EE}{\mathbb{E}}

% (a) (b) (c) labelling, without enumitem
\renewcommand{\theenumi}{\alph{enumi}}
\renewcommand{\labelenumi}{(\theenumi)}
\newenvironment{parts}
  {\begin{enumerate}\setlength{\itemsep}{0.45em}\setlength{\parskip}{0pt}}
  {\end{enumerate}}

\begin{document}

\begin{center}
  {\Large\bfseries HW 1 — Random Variables, Joints, Conditionals, Bayes}\\[0.4em]
  {\small Mathematical Foundations of Diffusion Models · Monsoon 2026 · IIT Jammu}
\end{center}

\medskip
\begin{center}
\begin{tabular}{@{}ll@{}}
\toprule
\textbf{Released} & Monday 10 August 2026 \\
\textbf{Due}      & Tuesday 18 August 2026, in class \\
\textbf{Marks}    & 50 points --- graded for completion \\
\bottomrule
\end{tabular}
\end{center}

\medskip
\noindent
\textbf{Graded for completion.} Credit is for a genuine attempt at every part, not
for arriving at the right answer. Attempt everything, and write down your
reasoning even where you are unsure — a wrong answer with visible working earns
full credit, a blank does not. Solutions are posted after the deadline; check your
own work against them, and bring anything that still does not make sense to class.

\medskip
\noindent
Show your working. Give exact fractions and decimals to three places.

\bigskip
\hrule
\bigskip

\section*{Problem 1}

Roll one fair six-sided die, once. Before the questions, here is what each symbol
in them means.

\begin{itemize}\setlength{\itemsep}{0.4em}
  \item $\Omega$ is the \textbf{sample space}: the set of every outcome the
        experiment could produce. For this die an outcome is just the face that
        comes up.
  \item $\omega$ is a single \textbf{outcome} — one particular thing that
        happened, one element of $\Omega$. Exactly one $\omega$ occurs each time
        you roll.
  \item An \textbf{event} is any subset $A \subseteq \Omega$: a yes/no question
        about the outcome. ``The face is even'' is the event $\{2,4,6\}$, and it
        occurs precisely when the $\omega$ that happened lies inside it.
  \item A \textbf{random variable} is a function $X : \Omega \to \mathbb{R}$ that
        attaches a number to \emph{each outcome}. It is not random and it is not a
        variable: it is an ordinary deterministic function. All the randomness sits
        in which $\omega$ turns up.
\end{itemize}

\noindent
Two random variables on this die:

\begin{itemize}\setlength{\itemsep}{0.4em}
  \item $X(\omega) = \omega$ — read the face and report it as a number. So
        $X(4) = 4$. (This one is the identity function; it looks like it is doing
        nothing, and that is the point: the outcome already \emph{is} a number
        here, which is not true in general.)
  \item $Y(\omega) = 1$ if $\omega$ is prime, and $0$ otherwise — report whether
        the face is prime. So $Y(3) = 1$ and $Y(4) = 0$.
\end{itemize}

\begin{parts}
  \item Write down $\Omega$ explicitly, and say how many distinct events this
        experiment has. Justify the count in one line. \pts{2}
  \item List the two sets $\{\omega : Y(\omega) = 1\}$ and
        $\{\omega : Y(\omega) = 0\}$. (Read the notation as ``all outcomes
        $\omega$ for which $Y$ returns that value''.) \pts{2}
  \item A classmate says ``$Y$ maps the event $\{2,3,5\}$ to $1$''. Explain why
        this is the wrong description of what $Y$ does, given what a random
        variable takes as its input. \pts{3}
  \item Compute $\PP(Y = 1)$ and $\PP(X \le 3,\ Y = 1)$. In each case name the
        event whose probability you are computing before you compute it. \pts{3}
\end{parts}

\total{10}

\bigskip
\hrule
\bigskip

\section*{Problem 2}

A workshop runs three machines, $A$, $B$ and $C$, which all make the same part.
Every finished part is inspected and given one of three quality grades:

\begin{center}
\begin{tabular}{cl}
\toprule
Grade & Meaning \\
\midrule
1 & within tight tolerance — ships as premium \\
2 & within tolerance — ships as standard \\
3 & out of tolerance — sent for rework \\
\bottomrule
\end{tabular}
\end{center}

\noindent
A day's output of 200 parts is inspected, and each part is recorded twice over:
which machine made it, and what grade it received. The counts are

\begin{center}
\begin{tabular}{lccc}
\toprule
        & grade 1 & grade 2 & grade 3 \\
\midrule
machine $A$   & 40 & 20 & 10 \\
machine $B$   & 15 & 35 & 20 \\
machine $C$   & 10 & 15 & 35 \\
\bottomrule
\end{tabular}
\end{center}

\noindent
Now pick one of the 200 parts uniformly at random, and let $X \in \{A,B,C\}$ be
the machine that made it and $Y \in \{1,2,3\}$ its grade. Both are random
variables on the same outcome — the part you picked.

The question behind the arithmetic is whether the machine a part came from tells
you anything about the grade it is likely to get. If it does, the workshop has a
machine worth investigating.

\begin{parts}
  \item Compute the row totals, column totals, and both marginal distributions
        $\PP(X = \cdot)$ and $\PP(Y = \cdot)$. \pts{2}
  \item Compute $\PP(Y = \cdot \mid X = A)$ and $\PP(Y = \cdot \mid X = C)$.
        Comparing the two, does knowing the machine change what you expect of the
        grade — that is, are $X$ and $Y$ dependent? \pts{3}
  \item A part is pulled off the standard-grade pile, so you know only that
        $Y = 2$. Compute $\PP(X = B \mid Y = 2)$, stating which total you divided
        by and why it is not 200. \pts{2}
  \item If machine and grade were unrelated, how many of the 200 parts would you
        expect to be grade-1 parts from machine $A$? Compare with the observed
        count, and say what the comparison suggests. \pts{3}
\end{parts}

\total{10}

\bigskip
\hrule
\bigskip

\section*{Problem 3}

A university holds a campus health camp and offers every student a quick
finger-prick blood-glucose screen for \textbf{type 2 diabetes}. Among students of
this age group roughly \textbf{1 in 100} has the condition, very often without
knowing it.

The screen is not a diagnosis; it only flags who should be sent for a proper
fasting-glucose or HbA1c test. Suppose it behaves as follows:

\begin{itemize}\setlength{\itemsep}{0.3em}
  \item it is positive for $95\%$ of students who \emph{do} have type 2 diabetes;
  \item it is also positive for $10\%$ of students who do \emph{not} — a single
        high reading can follow a recent meal, an illness, poor sleep, or ordinary
        day-to-day variation.
\end{itemize}

\noindent
Let $D$ be the event ``the student has type 2 diabetes'' and $+$ the event ``the
screen reads positive''. A student is picked at random from those screened.

\begin{parts}
  \item Write the three given numbers in conditional-probability notation. \pts{2}
  \item Compute $\PP(+)$: the fraction of screened students who get a positive
        reading. \pts{3}
  \item Compute $\PP(D \mid +)$: given a positive reading, the probability the
        student actually has the condition. \pts{3}
  \item The screen catches $95\%$ of true cases, yet your answer to (c) is far
        below $95\%$. Explain why. Then say which of the three given numbers you
        would change, and in which direction, to raise $\PP(D \mid +)$ the most.
        \pts{2}
\end{parts}

\noindent\textit{Check: redo (b) and (c) by imagining the camp screens 10{,}000
students, and filling in a $2 \times 2$ table of counts. Same answer, easier
arithmetic.}

\medskip
\noindent\textit{The figures above are illustrative and chosen to keep the
arithmetic clean; they are not the published operating characteristics of any
particular test. The structure, though, is real, and it is exactly why a positive
screening result triggers a confirmatory test rather than a diagnosis.}

\total{10}

\bigskip
\hrule
\bigskip

\section*{Problem 4}

Let $(X,Y)$ have joint density $f(x,y) = c\,(x+y)$ on $0 \le x \le 1$,
$0 \le y \le 2$, and $f = 0$ elsewhere.

\begin{parts}
  \item Find $c$. \pts{2}
  \item Find $f_X(x)$ and $f_Y(y)$ with their ranges, and verify each integrates
        to $1$. \pts{3}
  \item Find $f(y \mid x)$ and verify it integrates to $1$ over $y$. \pts{3}
  \item Are $X$ and $Y$ independent? Justify by explicit comparison. \pts{2}
\end{parts}

\total{10}

\bigskip
\hrule
\bigskip

\section*{Problem 5}

Two machines feed the same conveyor. Machine $1$ is the older unit and runs
$60\%$ of the output; machine $2$ runs the rest. At the end of the line every
item is inspected and its surface defects counted: $0$, $1$ or $2$.

Here is the difference from Problem 2, and it is the whole point. There, the
machine was written on the inspection sheet, so you observed \emph{both}
variables and could build the full table. Here the two machines' output mixes on
the shared conveyor, and by the time an item reaches inspection nobody knows
which machine made it. So the machine label $Z$ is \textbf{latent} — never
observed — while the defect count $X$ is all you get to see.

It still matters. If one machine is drifting out of adjustment you want to know
which one, so you can stop the line and service it; but the only evidence
available is defect counts. Recovering $Z$ from $X$ is an inference problem, and
Bayes' rule is what does it.

The two machines have different defect profiles:

\begin{center}
\begin{tabular}{lccc}
\toprule
                            & $X=0$ & $X=1$ & $X=2$ \\
\midrule
$\PP(X = \cdot \mid Z = 1)$ & 0.7 & 0.2 & 0.1 \\
$\PP(X = \cdot \mid Z = 2)$ & 0.2 & 0.3 & 0.5 \\
\bottomrule
\end{tabular}
\end{center}

\noindent
Pick an inspected item at random. Write $Z \in \{1,2\}$ for the machine that made
it and $X \in \{0,1,2\}$ for its defect count. You see $X$; you want $Z$.

\begin{parts}
  \item Using $p(z,x) = p(z)\,p(x \mid z)$, write the full $2 \times 3$ joint
        table and check it sums to $1$. \pts{2}
  \item Compute $\PP(X = x)$ for each $x$, stating which variable you summed
        out. This is the defect distribution the inspector actually sees, with the
        machine identity averaged away. \pts{2}
  \item An item is found with $x$ defects. Compute $\PP(Z = 1 \mid X = x)$ for
        each of $x = 0, 1, 2$ — the belief that the older machine made it. \pts{3}
  \item Compute $\EE[Z \mid X = x]$ for each $x$, and explain why these values
        need not lie in $\{1,2\}$ — and so what such a number is telling the
        engineer. \pts{2}
  \item Verify that $\EE\big[\EE[Z \mid X]\big] = \EE[Z]$. \pt{1}
\end{parts}

\noindent\textit{Parts (c) and (d) are, in miniature, what a diffusion model's
denoiser computes: a posterior over a hidden quantity given a noisy observation,
then its mean. The continuous version $\EE[x_0 \mid x_\sigma]$ appears in the
tutorial ``The World's Smallest Diffusion Model''.}

\total{10}

\end{document}
